\section{Computational evidence and its limits}
\label{exp:computation}
\label{exp:section}

The computations distinguish four questions: whether a formula reproduces
the implemented ray intersections, how strong a cut is on one corner,
what happens after LP reoptimization, and whether a rule improves a solver
run. A positive answer to an earlier question does not imply a positive
answer to a later one. We report retained experiments and corrected
records; no new optimization experiment was conducted for this paper.
The companion contains the compact records and their provenance.
\Cref{exp:protocols} gives cohort definitions and numerical conventions.

\subsection{Implementation fidelity and numerical corner values}
\label{exp:fidelity}

The source audit and instrumented runs concern SCIP~10.0.3. Across
610,193 ray comparisons from 7,945 ray-bearing attempts, the corrected
set model differs from SCIP's step on 1,211 rays (0.1985\%). The retained
classifications are 819 model-cutoff cases, 31 SCIP infinity classifications,
307 conservative bisection cases, and 54 rounding or root-formula cases.
The set case agrees in every attempt. The relative scaling of the
Case~4 eigen-coordinate blocks matters when
$\kappa_{\mathrm{SCIP}}\ne0$: an incorrect scaling need not describe
an $S$-free set. The bilinear experiments used
$\kappa_{\mathrm{SCIP}}=0$ and are unaffected by this formula pitfall.
They required a separate corner-oracle correction described below.

On the sampled MINLPLib corners, numerical degeneracy often leaves the
corner-bound criterion without a positive quantity to optimize. Of 4,578
ray-bearing attempts, 3,347 (73.1\%) have a numerically zero-rate face
meeting $S$. The classification treats
$\omega_j\le10^{-9}\max_i\omega_i$ as zero. Another 164 records have
all costs in the zero class. Reachability distinguishes zero corner value
from infinite corner value in the latter class.

The remaining 1,067 records, from 26 instances, have a positive value
under the reduced-space numerical test. Both $z_C$ and $\zK$ in the
reported ratios use costs floored at the same threshold. The median
SCIP-set ratio $z_C/\zK$ is 0.923 per corner
and 0.818 with equal instance weighting; 371 records (34.8\%) have ratio
below 0.5. These are numerical ratios. Support-at-most-two calculations,
generic three-ray KKT enumeration, and higher-dimensional solver estimates
have different guarantees; their counts and limitations appear in
\Cref{exp:corner-oracles}.

In 183 of the 371 low-ratio records, the determining numerically zero-rate
ray exits the SCIP set but has no detected one-ray intersection with $S$.
Four additional records have approximately $10^{-14}$ drift in the dumped
LP values. Their exact dumped-data ray restrictions reach $S$ only after
steps of order $10^{18}$. Whether this drift occurs at the true LP corner
is undetermined. Those four records do not establish a geometric gap.

\begin{table}[t]
\centering
\caption{Implementation and rejection diagnostics. Each row has its own
denominator. Rescaling passes concern the first restriction piece only.}
\label{exp:diagnostic-table}
\begin{tabular}{@{}lrr@{}}
\toprule
Observation & Count & Fraction \\
\midrule
Classified step discrepancies & $1,211/610,193$ & $0.20\%$ \\
Numerical zero-rate face meets $S$ & $3,347/4,578$ & $73.1\%$ \\
Generated cuts, all MINLPLib attempts & $29,022/76,914$ & $37.7\%$ \\
Dynamism aborts, all attempts & $38,105/76,914$ & $49.5\%$ \\
Free-basis-status aborts, all attempts & $9,784/76,914$ & $12.7\%$ \\
First-piece rescaling passes, sample & $1,728/1,799$ & $96.1\%$ \\
First-piece rescaling passes, full diagnostic & $31,926/34,834$ & $91.7\%$ \\
\bottomrule
\end{tabular}
\end{table}

Consistent ray rescaling changes the intersection parameter inversely
and preserves the mathematical cut after conversion to the original
coordinates. The inspected implementation does not normalize these rays.
Rescaling makes many recorded first-piece dynamism tests pass
(\Cref{exp:diagnostic-table}). This is a diagnostic of a noninvariant
rejection test. It does not count valid cuts recovered by a complete
implementation: later pieces and numerical safeguards remain untested.
The observed LP-seen rates are lower bounds, since a cut may enter and
leave between snapshots. The attempt counts describe the instrumented
binary, including its overhead on time-limited runs.

\subsection{A corner optimum and an LP re-solve are different quantities}
\label{exp:one-cut}

The one-cut experiment starts from a McCormick LP, selects its most
violated bilinear term, and compares cuts after re-solving the full LP.
Its denominator is the gap between the initial LP and a global solve
enforcing that one term. The retained sample has 47 small and 73 larger
corners. Ten old corner values were infeasible underestimates caused by
antiparallel projected rays. Rational lower and upper checks correct those
ten values; \Cref{exp:ten-corrections} identifies them.

\begin{table}[t]
\centering
\caption{Corrected one-cut experiment: mean fraction of the single-term
gap. The orbit search attains the numerical corner value to bisection
resolution on these 120 corners.}
\label{exp:one-cut-table}
\begin{tabular}{@{}p{0.47\textwidth}cc@{}}
\toprule
Quantity & \shortstack{4 variables, 4 products\\($n=47$)} & \shortstack{6 variables, 8 products\\($n=73$)} \\
\midrule
SCIP-set corner increment & 0.724 & 0.616 \\
Orbit/corner increment & 0.783 & 0.675 \\
SCIP-cut LP re-solve & 0.895 & 0.870 \\
Orbit-cut LP re-solve & 0.931 & 0.910 \\
Objective-parallel corner cut, LP re-solve & 0.783 & 0.675 \\
\bottomrule
\end{tabular}
\end{table}

The objective-parallel cut gains exactly the corner increment. Tilted
cuts can gain more through other LP rows. The orbit cut improves on the
SCIP-model cut by more than 0.01 of the single-term gap on 62 corners,
and loses by more than 0.01 on 22. Thus even numerical corner attainment
does not give LP dominance. The numerical minimum orbit/corner ratio is
about $0.999999989$; this is not a collection of 120 exact certificates.
SCIP's ratio is below 0.9 on 27 corners and below 0.5 on three.

The minor experiment leads to the same distinction for signature $(2,2)$.
It contains 86 and 71 LP corners from random $3\times3$ and $4\times4$
models. On all 157, the orbit search attains the numerical corner value,
whereas SCIP's set does so on 12. Yet the maximum-margin choice decreases
mean full-LP gap closure relative to SCIP by 0.0758 and 0.0651.
The orbit-optimal set nearest SCIP has mean differences $+0.0225$ and
$+0.0057$; the comparisons do not establish a mean advantage.
These are single-minor LP re-solves, without branch-and-bound integration.
The exact minor examples in \Cref{mi:examples} establish geometric
obstructions independently of these frequencies.

\subsection{Later rounds and explored policies}
\label{exp:multiround}

The larger experiment is an external root cutting loop. At each round
it re-solves the LP, forms its corner, and cuts violated bilinear terms.
There are 220 paired constructed instances: 60 each at sizes $4\times4$
and $6\times8$, and 50 each at $8\times12$ and $10\times20$.
Here the denominator is the gap to a global solve enforcing all bilinear
equations. The implemented orbit search floors reduced costs and uses
floating-point bisection. We call it the orbit rule without claiming exact
optimization for raw zero reduced costs.

\begin{table}[t]
\centering
\caption{Mean root-gap fraction closed by the external loop. Intervals
are 95\% paired-$t$ intervals for orbit minus SCIP model after round 10.}
\label{exp:multiround-table}
\small
\begin{tabular}{@{}lcccc@{}}
\toprule
Size ($n$) & SCIP r1 & Orbit r1 & SCIP/orbit r10 & Difference r10 \\
\midrule
$4\times4$ (60) & 0.757 & 0.761 & 0.974 / 0.966 & $-0.008\ [-0.016,+0.001]$ \\
$6\times8$ (60) & 0.754 & 0.755 & 0.968 / 0.957 & $-0.011\ [-0.019,-0.002]$ \\
$8\times12$ (50) & 0.804 & 0.799 & 0.974 / 0.969 & $-0.005\ [-0.010,+0.000]$ \\
$10\times20$ (50) & 0.705 & 0.699 & 0.937 / 0.912 & $-0.025\ [-0.045,-0.005]$ \\
\bottomrule
\end{tabular}
\end{table}

Let $\Delta_r$ denote the pooled orbit-minus-SCIP difference after round $r$.
The estimates and 95\% paired intervals are
\begin{align*}
\Delta_1&=-0.001, &[-0.016,+0.014],\\
\Delta_{10}\ \text{and}\ \Delta_{20}&=-0.012, &[-0.017,-0.006].
\end{align*}
The antiparallel-ray correction changes the first-round orbit result by
$+0.011\ [+0.004,+0.019]$ against the bug-affected implementation, but
the round-10 difference is $-0.001\ [-0.003,+0.002]$.
It therefore does not explain the later-round reversal.

Switching after one orbit root round to SCIP-model rounds already leaves
a pooled round-10 deficit $-0.006\ [-0.011,-0.002]$. Cross-evaluations
of retained states after early schedules report different amounts of
subsequent progress under the tested rules. Bound attainment and ties
alone do not explain the loss:
nonattaining, low-tie policies also lose. Shorter steps on many rays are
a plausible common feature. The intervention contrasting a majority-ray
selection rule with random orbit selection is inconclusive and does not
hold orbit exposure constant. These observations do not identify a causal
mechanism.

The efficacy policy chooses, per term, the largest Euclidean efficacy
among the SCIP cut, the orbit cut, two cost-perturbed orbit choices, and
a Euclidean-depth orbit choice. Selected after examining the records,
it has pooled
advantage 0.004 in AUC, defined as the mean gap fraction at rounds 1--10,
with Holm-adjusted $p=0.025$. The Holm family has 14 complete-cohort rule
comparisons against SCIP and is applied separately to AUC and round 10;
the latter adjusted $p$ is 0.33. Comparisons of switching from orbit to
SCIP-model rounds against continuing orbit also depend on the tested
family and statistical test (\Cref{exp:policy-details}). Both findings need fresh
data. The convergence hypotheses of the mathematical results are not
verified by these finite trajectories.

\subsection{Selection inside SCIP}
\label{exp:selection}

The inside-SCIP experiment searches alternative admissible directions
using 24 half-circle samples with local refinement in two dimensions,
and projected subgradient ascent in higher dimensions. This is a
heuristic search. Root-gap closure is measured against recorded reference
values on complete, comparable cases; its denominator differs from the
one-term experiment above.

\begin{table}[t]
\centering
\caption{Paired mean changes in root-gap fraction inside SCIP~10.0.3.
The comparisons within each seed use the same eligible cohort.}
\label{exp:root-table}
\begin{tabular}{@{}lccc@{}}
\toprule
Comparison & Seed 0 ($n=251$) & Seed 1 ($n=250$) & Seed 2 ($n=248$) \\
\midrule
Point rule minus disabled cuts & $+0.114$ & $+0.104$ & $+0.106$ \\
Corner criterion minus point rule & $-0.008$ & $-0.007$ & $-0.010$ \\
Efficacy criterion minus point rule & $+0.001$ & $+0.003$ & $-0.004$ \\
\bottomrule
\end{tabular}
\end{table}

The corner criterion provides no consistent improvement in this sample.
In a full summary scan of 305 instances and 109,389 corners, 81.6\% of
the floored criterion minima occur at numerically zero-rate rays. The
positive surrogate then measures a flooring convention rather than a
positive true corner gain. Degeneracy alone does not explain every loss.
The raw individual corners from this full scan were not retained, so
the companion preserves its summaries rather than a full reconstruction.

Full solves do not repair the inference. Row capture in the patched
build changes some tree paths; the subsequent 120 stock reruns were made
in a different batch. Their CPU timing comparison is uncontrolled and
is excluded from the performance conclusions. The 300 debug-solution
runs are a separate validity audit with incomplete objective mappings
and unresolved diagnostics, not ordinary full-solve benchmarks or a
clean validity certificate (\Cref{exp:debug}). The evidence establishes
neither a stock-SCIP timing benefit nor a generally better selection rule.
